\section{Audit Relations and Guarantees}
\label{sec:requirements}

Section~\ref{sec:problem} defines the completed episode and its
audit interfaces. We now formalize the three audit relations
summarized in Figure~\ref{fig:teaser}. Evidence authorization
establishes the data view permitted by the agent's declaration.
Execution reproducibility determines whether the final program
reproduces the stated answer within that view. Semantic correctness
independently evaluates whether the program implements the requested
task. The first two relations determine the execution-provenance
status; the third determines the semantic status.

\subsection{Evidence authorization}
\label{sec:authorization}

The declared evidence first determines the data view against which
the recorded execution is audited.

\begin{definition}[Evidence authorization]
\label{def:authorization}
A completed episode $E$ is evidence-authorized when every declared
evidence identifier resolves and
\[
F_{\mathrm{host}}\subseteq F_{\mathrm{res}},
\qquad
F_{\mathrm{res}}=R(F_{\mathrm{decl}}).
\]
\end{definition}

The two conditions produce distinct audit evidence. An unresolved
declaration identifies a claimed source that cannot be included in
the evidence view. An undeclared host read identifies an observed
execution dependency outside that view. Canonical resolution
precedes the subset test so that aliases and path variants refer to
the same data object.

Evidence authorization establishes $F_{\mathrm{res}}$ as the input
view for restricted replay. It does not yet establish that the final
program reproduces the stated answer.The running example violates this relation because 16
fraud-and-other files read by $P^*$ are absent from
$F_{\mathrm{res}}$.

\subsection{Execution reproducibility}
\label{sec:reproducible}

Let
\[
\rho
=
\operatorname{Exec}_{\mathcal R_{\mathrm{replay}}}
(P^*,F_{\mathrm{res}})
=
(s_{\mathrm{replay}},
 v_{\mathrm{replay}},
 F_{\mathrm{replay}},
 \operands_{\mathrm{replay}})
\]
be the restricted replay record defined in
Section~\ref{sec:statuses}. A replay certificate is issued exactly
when
\begin{equation}
\begin{aligned}
\operatorname{Cert}_{\mathrm{replay}}(E)=1
\ \Longleftrightarrow\quad
&F_{\mathrm{res}}=R(F_{\mathrm{decl}})
  \text{ is complete},\\
&F_{\mathrm{host}}\subseteq F_{\mathrm{res}},
  \quad s_{\mathrm{host}}=\mathrm{success},\\
&\rho=
  \operatorname{Exec}_{\mathcal R_{\mathrm{replay}}}
  (P^*,F_{\mathrm{res}}),\\
&s_{\mathrm{replay}}=\mathrm{success},
  \quad F_{\mathrm{replay}}\subseteq F_{\mathrm{res}},\\
&\operatorname{cmp}_{\tol}
  (v_{\mathrm{replay}},\hat a)=1.
\end{aligned}
\label{eq:replaycert}
\end{equation}

The replay witness stores the final program snapshot, resolved
evidence identifiers, host and replay reads, replay output,
comparator, and tolerance.

\begin{theorem}[Relative soundness of replay]
\label{thm:relative}
Under the trusted execution model of
Section~\ref{sec:threat}, a replay certificate establishes that
$P^*$ reproduces $\hat a$ within the resolved declared evidence
view, and that all observed host and replay data reads remain
inside that view.
\end{theorem}

\noindent\emph{Proof sketch.}
Equation~\ref{eq:replaycert} requires complete evidence resolution,
authorized host reads, restricted replay, authorized replay reads,
and output agreement. The trusted logger, resolver, replay runtime,
and comparator make these conditions observable and enforceable.
The full proof appears in Appendix~\ref{app:proofs}.

The guarantee is relative to the recorded final program and the
resolved evidence view. It establishes how the stated answer was
reproduced, while leaving the question--program relation to the
semantic audit. Its restricted replay additionally returns a null result, so
$\operatorname{cmp}_{\tol}(v_{\mathrm{replay}},\hat a)=\bot$.

\subsection{Semantic correctness}
\label{sec:nonimplication}

Semantic correctness evaluates whether $P^*$ implements the
task-specific specification $V_q$ for question $q$:
\[
S(E;V_q)
\in
\{
\mathrm{correct},
\mathrm{incorrect},
\mathrm{unknown}
\}.
\]
Depending on the task, this relation may inspect the requested
aggregation, predicates, selected scope, output type, or other
executable constraints. It is evaluated independently of whether
the recorded program reproduces the stated answer.

\begin{proposition}[Semantic non-implication]
\label{prop:semantic}
There exist completed episodes for which
\[
\operatorname{Cert}_{\mathrm{replay}}(E)=1
\qquad\text{and}\qquad
S(E;V_q)=\mathrm{incorrect}.
\]
Consequently, execution-provenance consistency does not imply
semantic correctness.
\end{proposition}

\noindent\emph{Proof sketch.}
Choose two programs with different task semantics but identical
outputs on the current data. Replay certifies the submitted
program's output, while the task-specific validator rejects its
operation. Appendix~\ref{app:proofs} gives the full argument.

For example, a sum and an average coincide after filtering to one
row, and an incorrect predicate may leave the result unchanged on
the current database. Such episodes are reproducible but violate
the task specification. The same separation applies to the bounded
closure introduced in Section~\ref{sec:closure}: derivability from
reached operands does not identify the operation requested by the
question. This distinction motivates the independent semantic
branch in Figure~\ref{fig:teaser} and the same-output SQL evaluation
in Section~\ref{sec:rq3}.